\section{Pewne definicje związane z algebrami.}

Wprowadzimy teraz definicje,
	które dzięki pojęciu algebry będą miały sens dla wszystkich struktur,
	które są modelami algebr.

\subsection{Podalgebry.}

\begin{definition}
    \textbf{Podmodelem modelu $M$ algebry $(\Sigma, E)$} nazywamy podzbiór $N \subseteq M$,
		który jest modelem algebry $(\Sigma, E)$ z tymi samymi operacjami co $M$.
\end{definition}

\begin{example}
	$2\mathbb{N}$, czyli liczby parzyste naturalne stanowią podpółgrupę liczb naturalnych z dodawaniem.
\end{example}
\begin{example}
	$\mathbb{Q} \setminus \{0\}$ stanowi podgrupę przemienną grupy przemiennej $\mathbb{R} \setminus \{0\}$ z mnożeniem.
\end{example}

\subsection{Homomorfizmy.}

\begin{definition}
    Niech $\mathfrak{A} = (\Sigma, E)$ będzie algebrą.
	\textbf{Homomorfizmem} (lub \textbf{morfizmem}) modeli $\mathfrak{A}$ nazywamy dowolną funkcję
		$h: A \to B$, gdzie $A, B$ są dowolnymi modelami $\mathfrak{A}$,
		i dla każdego symbolu $f^n \in \Sigma$ o arności $n$,
			któremu w $A$ odpowiada funkcja $f_A: A^n \to A$,
			a w $B$ funkcja $f_B: B^n \to B$,
		 zachodzi
	\begin{equation}
		\forall_{x_1, \dots, x_n \in A} \; h(f_A(x_1, \dots, x_n)) = f_B(h(x_1), \dots, h(x_n)).
	\end{equation}
\end{definition}

\begin{example}
	Funkcja
	\begin{align*}
		&f: \mathbb{Q} \to \mathbb{R},\\
		&f(x) = 5x,
	\end{align*}
	jest homomorfizmem grup $\mathbb{Q}$ i $\mathbb{R}$ z dodawaniem.
	Rzeczywiście, dla każdych $x, y \in \mathbb{Q}$ mamy
	\begin{align}
		5 \cdot (x + y) &= 5x + 5y, \\
		5 \cdot (-x) &= -(5 \cdot x), \\
		5 \cdot 0 &= 5 \cdot 0.
	\end{align}
\end{example}

\begin{example}
	Funkcja
	\begin{align*}
		&f: \mathbb{R}_+ \to \mathbb{R}, \\
		&f(x) = \log(x)
	\end{align*}
	jest homomorfizmem grup $\mathbb{R}_+$ z mnożeniem i $\mathbb{R}$ z dodawaniem.
	Rzeczywiście, dla każdych $x, y \in \mathbb{R}_+$ mamy
	\begin{align}
		\log(x \cdot y) &= \log(x) + \log(y), \\
		\log(x^{-1}) &= -\log(x), \\
		\log(1) &= 0.
	\end{align}
\end{example}

Z pewnych niewyjaśnionych powodów,
    w przypadku homomorfizmów nie używa się pojęć takich jak iniekcja, bijekcja i surjekcja.
Zamiast tych pochodzących z francuskiego wyrazów mamy greckie odpowiedniki.

\begin{definition}
    Homomorfizm nazywamy
    \begin{itemize}
        \item \textbf{epimorfizmem}, jeśli jest surjekcją,
        \item \textbf{monomorfizmem}, jeśli jest iniekcją,
        \item \textbf{izomorfizmem}, jeśli jest bijekcją,
        \item \textbf{endomorfizmem}, jeśli jego dziedzina i kodziedzina są równe, t.j. jest postaci $h:A \to A$,
        \item \textbf{automorfizmem}, jeśli jest endomorfizmem i izomorfizmem.
    \end{itemize}
\end{definition}

\begin{lemma}
	Niech $\mathfrak{A} = (\Sigma, E)$ będzie algebrą.
	Niech $A, B, C$ będą modelami $\mathfrak{A}$,
		a $h: A \to B$, $g: B \to C$ homomorfizmami tych modeli.
	Wtedy $g \circ h: A \to C$ jest homomorfizmem.
\end{lemma}
\begin{proof}
	Niech $f^n \in \Sigma$ będzie $n$-arnym symbolem,
		a $f_A, f_B, f_C$ odpowiadającymi mu funkcjami w odpowiednio $A, B, C$.
	Mamy teraz dla dowolnych $x_1, \dots, x_n \in A$
	\begin{align}
		(g \circ h)(f_A(x_1, \dots, x_n)) &= \\
			&= g(h(f_A(x_1, \dots, x_n)) \\
			&= g(f_B(h(x_1), \dots, h(x_n))) \\
			&= f_C(g(h(x_1)), \dots, g(h(x_n)))) \\
			&= f_C((g \circ h)(x_1), \dots, (g \circ h)(x_n)))),
\end{align}
	co pokazuje, że $(g \circ h)$ jest homomorfizmem.
\end{proof}

\begin{lemma}
	Jeśli $h: A \to B$ jest izomorfizmem dwóch modeli algebry $\mathfrak{A} = (\Sigma, E)$,
		to funkcja odwrotna $h^{-1}: B \to A$ istnieje i też jest izomorfizmem.
\end{lemma}
\begin{proof}
	Istnienie $h^{-1}$ i to,
		że jest bijekcją,
		jest oczywiste z faktu,
		że $h$ jest bijekcją.
	Wystarczy więc pokazać,
		że $h^{-1}$ jest homomorfizmem.
	Dla dowolnego symbolu $f^n \in \Sigma$ reprezentowanego przez funkcje
		$f_A: A^n \to A$ i $f_B: B^n \to B$,
		i dowolnych $x_1, \dots, x_n \in B$ mamy
	\begin{align}
		h(f_A(h^{-1}(x_1), \dots, h^{-1}(x_n))) &= \\
			&= f_B(h(h^{-1}(x_1)), \dots, h(h^{-1}(x_n))) \\
			&= f_B(x_1, \dots, x_n) \\
			&= h(h^{-1}(f_B(x_1, \dots, x_n))).
	\end{align}
	Ale ponieważ $h$ jest iniekcją, to
	\begin{equation}
		f_A(h^{-1}(x_1), \dots, h^{-1}(x_n)) = h^{-1}(f_B(x_1, \dots, x_n)),
	\end{equation}
	czyli $h^{-1}$ jest homomorfizmem.
\end{proof}

\begin{definition}
	Mówimy, że modele $A, B$ algebry $(\Sigma, E)$ są \textbf{izomorficzne},
		jeśli istnieje izomorfizm $f: A \to B$.
	Fakt ten zapisujemy $A \cong_f B$ lub po prostu $A \cong B$.
\end{definition}
Jeśli dwa modele są izomorficzne, to w pewnym sensie są identyczne---różnią się tylko formą zapisu.
Poza tym wszystkie algebraiczne własności,
	które zachodzą w jednym modelu,
	zachodzą też w drugim.
W niektórych dziedzinach matematyki izomorficzne modele identyfikuje się tak bardzo,
	że zamiast używać znaku $A \cong B$ używa się zwykłej równości $=$.
My jednak zostaniemy przy $\cong$.

\begin{example}
	$\mathbb{R}$ z dodawaniem jest grupą abelową izomorficzną z grupą abelową $\mathbb{R}^{1 \times 1}$ z dodawaniem macierzy.
	Izomorfizmem jest tu funkcja
	\begin{align*}
		&f: \mathbb{R} \to \mathbb{R}^{1 \times 1},\\
		&f(x) = [x].
	\end{align*}
\end{example}

\begin{example}
$\mathbb{Z}$ z dodawaniem (ozn. $\mathbb{Z}^+$)
	jako monoid nie jest izomorficzne z $\mathbb{Z}$ z mnożeniem (ozn. $\mathbb{Z}^*$)
	jako monoid.
\end{example}
\begin{proof}
Przypuśćmy, że istnieje izomorfizm $f: \mathbb{Z}^+ \to \mathbb{Z}^*$. Istnieje więc też odwrotność tego izomorfizmu $f^{-1}$, która też jest izomorfizmem. Mamy zatem
\begin{equation}
    f^{-1}(1) = 0 = 1 + (-1) = f^{-1}(f(1)) + f^{-1}(f(-1)) = f^{-1}(f(1) \cdot f(-1)).
\end{equation}
Skoro $f^{-1}$ jest izomorfizmem, to 
\begin{equation}
    1 = f(1) \cdot f(-1), 
\end{equation}
więc
\begin{equation}
    f(1) = 1 = f(-1) \text{ lub } f(1) = -1 = f(-1).
\end{equation}
Wtedy jednak korzystając z tego, że $f$ jest izomorfizmem dostajemy $1 = -1$, co jest sprzecznością.
\end{proof}

\subsection{Wolne algebry.}

\begin{definition}
    Wolnym modelem algebry $(\Sigma, E)$ generowanym przez zbiór $G$ (tak zwany \textbf{zbiór generatorów}),
    nazywamy model algebry $(\Sigma, E)$, którego nośnikiem jest zbiór wszystkich wyrażeń,
		które utworzyć możemy z elementów $\Sigma$ i $G$, traktują elementy $\Sigma$ jak
		funkcje z odpowiednimi arnościami.
	Dodatkowo,
		dwa wyrażenia uznajemy za równe,
		jeśli jedno możemy przekształcić w drugie używając równań z $E$.
\end{definition}
Mówimy też, że model jest \textit{nad zbiorem $G$}.
Znowu nie jest to w pełni formalna definicja,
	jednak bardzo intuicyjnie oddająca istotę wolnych modeli.
Definicja w pełni formalna wymagałaby albo wprowadzenia pewnych formalnych pojęć,
	których nie będziemy wprowadzać (np. ilorazy modeli przez podmodele),
	albo byłaby bardzo abstrakcyjna (pojęcie własności uniwersalnych).

Zwykle jako zbiór $G$ bierze się pewne formalne symbole, czyli jakieś litery.

\begin{example}
	Grupą wolną na zbiorze $\{a\}$ nazywamy grupę, której elementami są
	$e, a, a^{-1}, a^2, a^{-2}, \dots$.
	Zachodzą w niej m.in. równości
	\begin{align}
		aa^{-1} &= e \\
		eae &= a \\
		a^{-1}a &= e
	\end{align}
\end{example}

W pewnym sensie grupą wolną na zbiorze $G$ jest najprostrza grupa,
	która zawiera całe $G$ jako swoje elementy.
Najprostsza,
	bo nie zachodzą w niej inne równania,
	niż te,
	które stanowią aksjomaty grupy,
	oraz te,
	które można z nich wyprowadzić.
Podobnie jest dla innych struktur algebraicznych.

\begin{example}
	$\mathbb{Z}$ wraz z dodawaniem stanowi wolną grupę przemienną nad zbiorem $\{1\}$.
\end{example}

\begin{example}
	Zbiór skończonych ciągów liczb z $\mathbb{N}$ z konkatenacją stanowi wolny monoid nad $\mathbb{N}$.
\end{example}

\begin{example}
	$\mathbb{Z}$ wraz z mnożeniem nie stanowi wolnej grupy przemiennej,
		bo dla dowolnego elementu $x$ mamy
		\begin{equation}
		0x = 0
		\end{equation}
		czego nie możemy wyprowadzić z aksjomatów grupy przemiennej.
	Wystarczy podać grupę przemienną,
		w której nie ma elementu,
		który działa tak jak $0$ w $\mathbb{Z}$ z mnożeniem.
	Jest nią przykładowo $\mathbb{Z}$ z dodawaniem.
\end{example}

\subsection{Zadania.}

\begin{exercise}
	Zapisz definicję homomorfizmu monoidów nie używając pojęć algebry uniwersalnej.
\end{exercise}

\begin{exercise}
	Zapisz definicję homomorfizmu grup nie używając pojęć algebry uniwersalnej.
\end{exercise}

\begin{exercise}
	Sprawdź, czy następujące funkcje
	\begin{enumerate}
		\item $f(x) = 2x$,
		\item $g(x) = 2 + x$
		\item $h(x) = x^2$,
	\end{enumerate}
	są endomorfizmami w
	\begin{enumerate}
		\item monoidzie $\mathbb{Z}$ z mnożeniem,
		\item grupie $\mathbb{Z}$ z dodawaniem.
	\end{enumerate}
\end{exercise}

\begin{exercise}
	Niech $(A, +), (B, \cdot)$ będą dowolnymi grupami.
	Udowodnij,
		że jeśli $f: A \to B$ jest funkcją spełniającą
	\begin{equation}
		f(x + y) = f(x) \cdot f(y)
	\end{equation}
	dla każdych $x, y \in A$, to $f$ jest homomorfizmem grup.
\end{exercise}

\begin{exercise}
	Udowodnij,
		że jeśli $\mathfrak{A}$ jest algebrą,
		a $\mathcal{A}$ zbiorem modeli $\mathfrak{A}$,
		to relacja $\cong$ na elementach $\mathcal{A}$ jest relacją równoważności.
\end{exercise}

\begin{exercise}
	Niech $A$ będzie grupą abelową, a $B$ dowolną grupą.
	Udowodnij, że jeśli $f: A \to B$ jest epimorfizmem,
		to $B$ jest grupą abelową.
\end{exercise}

\begin{exercise}
	Udowodnij, że jeśli $A, B$ są zbiorami i $|A| = |B|$,
		to grupa wolna nad $A$ i grupa wolna nad $B$ są izomorficzne.
\end{exercise}

